\paragraph{Poisson--Cauchy 粗粒化的六阶信息精化：反向熵、零点结构与重参数化猜想}
在命题 \ref{con:xi-poisson-coarsegraining-second-order-normal-form} 与命题
\ref{prop:xi-poisson-kl-dissipation-fractional-fisher}
给出的二阶正规形与耗散恒等式框架内，本段把六阶层的可见信息进一步拆分为
正向/反向相对熵系数、一般 \(f\)-散度曲率项、\(L^p\) 尺度常数、差分剖面零点刚性，
以及与 Jensen 缺陷证书链的重参数化对齐。

\begin{theorem}[\texorpdfstring{$D_{\mathrm{KL}}$}{KL} 的 \texorpdfstring{$t^{-6}$}{t^-6} 精确二项展开：三、四阶中心矩入系数]\label{thm:xi-poisson-kl-sixth-order-refined-binomials}
\leanverified{paper\_xi\_poisson\_kl\_sixth\_order\_refined\_binomials}
设 \(\tilde\nu\) 为 \(\RR\) 上概率测度，记
\[
m_k:=\int_{\RR}(\gamma-\bar\gamma)^k\,\dd\tilde\nu(\gamma),\qquad
\sigma^2:=m_2>0,\qquad
\bar\gamma:=\int_{\RR}\gamma\,\dd\tilde\nu(\gamma),
\]
并假设
\(
\int_{\RR}\abs{\gamma-\bar\gamma}^5\,\dd\tilde\nu(\gamma)<\infty
\)。
定义
\[
P_t(x):=\frac{1}{\pi}\frac{t}{x^2+t^2},\qquad
\tilde h_t:=P_t*\tilde\nu,\qquad
\tilde g_t(x):=P_t(x-\bar\gamma).
\]
则当 \(t\to\infty\) 时，
\begin{equation}\label{eq:xi-poisson-kl-sixth-order-refined-binomials}
D_{\mathrm{KL}}(\tilde h_t\mid \tilde g_t)
=\frac{\sigma^4}{8t^4}
+\frac{1}{t^6}\left(
\frac{\sigma^6}{64}
+\frac{3m_3^2}{32}
-\frac{\sigma^2m_4}{8}
\right)
+o(t^{-6}).
\end{equation}
\end{theorem}

\begin{proof}
令 \(Z:=\gamma-\bar\gamma\)。对 \(x\mapsto \tilde g_t(x)=P_t(x-\bar\gamma)\) 作四阶 Taylor 展开并取积分余项，得到
\[
\tilde h_t
=\tilde g_t+\frac{m_2}{2}\tilde g_t''
-\frac{m_3}{6}\tilde g_t'''
+\frac{m_4}{24}\tilde g_t^{(4)}+R_t.
\]
由 \(\|P_t^{(5)}\|_\infty\asymp t^{-6}\) 及 \(\EE|Z|^5<\infty\)，有 \(\|R_t\|_\infty=O(t^{-6})\)。
记
\[
\delta_t(x):=\frac{\tilde h_t(x)}{\tilde g_t(x)}-1,\qquad
x=\bar\gamma+ty,\qquad
G(y):=\frac{1}{\pi(1+y^2)}.
\]
则
\[
\delta_t(\bar\gamma+ty)
=\frac{\sigma^2}{t^2}U_2(y)
+\frac{m_3}{t^3}V_3(y)
+\frac{m_4}{t^4}W_4(y)
+O(t^{-5}),
\]
其中
\[
U_2(y)=\frac{3y^2-1}{(1+y^2)^2},\quad
V_3(y)=\frac{4y(y^2-1)}{(1+y^2)^3},\quad
W_4(y)=\frac{5y^4-10y^2+1}{(1+y^2)^4}.
\]
由 \(\int \tilde g_t\delta_t=0\) 与 \(\delta_t=O(t^{-2})\) 得
\[
D_{\mathrm{KL}}(\tilde h_t\mid \tilde g_t)
=\frac12\int_{\RR}G\,\delta_t^2\,\dd y
-\frac16\int_{\RR}G\,\delta_t^3\,\dd y
+o(t^{-6}).
\]
利用 \(U_2,W_4\) 为偶、\(V_3\) 为奇，得到 \(t^{-5}\) 交叉项消失。再用 Cauchy 权重积分常数
\[
\int_{\RR}G\,U_2^2\,\dd y=\frac14,\qquad
\int_{\RR}G\,V_3^2\,\dd y=\frac{3}{16},
\]
\[
\int_{\RR}G\,U_2W_4\,\dd y=-\frac18,\qquad
\int_{\RR}G\,U_2^3\,\dd y=-\frac{3}{32},
\]
即可得到 \eqref{eq:xi-poisson-kl-sixth-order-refined-binomials}。
\end{proof}

\begin{proposition}[一般 \texorpdfstring{$f$}{f}-散度的 \texorpdfstring{$t^{-6}$}{t^-6} 二项展开]\label{prop:xi-poisson-fdiv-sixth-order-refined-binomials}
在定理 \ref{thm:xi-poisson-kl-sixth-order-refined-binomials} 的矩条件下，
设 \(f\) 为凸函数，满足 \(f(1)=0\)，并在 \(1\) 的邻域 \(C^3\)。定义
\[
D_f(\tilde h_t\mid \tilde g_t)
:=\int_{\RR}\tilde g_t(x)\,
f\!\left(\frac{\tilde h_t(x)}{\tilde g_t(x)}\right)\dd x.
\]
则当 \(t\to\infty\) 时，
\begin{equation}\label{eq:xi-poisson-fdiv-sixth-order-refined-binomials}
D_f(\tilde h_t\mid \tilde g_t)
=\frac{f''(1)\sigma^4}{8t^4}
+\frac{1}{t^6}\left[
f''(1)\!\left(\frac{3m_3^2}{32}-\frac{\sigma^2m_4}{8}\right)
-\frac{f'''(1)\sigma^6}{64}
\right]
+o(t^{-6}).
\end{equation}
\end{proposition}

\begin{proof}
在 \(u=0\) 处展开
\[
f(1+u)=f(1)+f'(1)u+\frac{f''(1)}{2}u^2+\frac{f'''(1)}{6}u^3+O(u^4).
\]
代入 \(u=\delta_t\) 后，线性项由 \(\int\tilde g_t\delta_t=0\) 消失，且 \(\delta_t=O(t^{-2})\) 使四次及以上项积分为 \(O(t^{-8})\)。于是
\[
D_f(\tilde h_t\mid \tilde g_t)
=\frac{f''(1)}{2}\int_{\RR}G\,\delta_t^2\,\dd y
+\frac{f'''(1)}{6}\int_{\RR}G\,\delta_t^3\,\dd y
+o(t^{-6}).
\]
再代入定理 \ref{thm:xi-poisson-kl-sixth-order-refined-binomials} 证明中的同一 \(\delta_t\) 展开和积分常数即可。
\end{proof}

\begin{theorem}[Poisson 差分剖面的 \texorpdfstring{$L^p$}{Lp} 普适尺度律]\label{thm:xi-poisson-lp-universal-scaling-refined}
在命题 \ref{con:xi-poisson-coarsegraining-second-order-normal-form} 的一致二阶展开与误差控制下，
对任意 \(p\in(1,\infty]\) 成立
\begin{equation}\label{eq:xi-poisson-lp-universal-scaling-refined}
\lim_{t\to\infty}\frac{t^{3-1/p}}{\sigma^2}\,\|\tilde h_t-\tilde g_t\|_{L^p(\RR)}=C_p,
\end{equation}
其中
\[
C_p:=\left(\int_{\RR}\left|
\frac{1}{\pi}\frac{3y^2-1}{(1+y^2)^3}
\right|^p\dd y\right)^{1/p}\quad(p<\infty),
\]
\[
C_\infty:=\sup_{y\in\RR}\left|
\frac{1}{\pi}\frac{3y^2-1}{(1+y^2)^3}
\right|=\frac{1}{\pi}.
\]
特别地，
\begin{equation}\label{eq:xi-poisson-lp-c2-refined}
C_2=\sqrt{\frac{3}{16\pi}}.
\end{equation}
\end{theorem}

\begin{proof}
由命题 \ref{con:xi-poisson-coarsegraining-second-order-normal-form}，
\[
\tilde h_t-\tilde g_t=\sigma^2 r_t+\varepsilon_t,
\]
且
\[
r_t(\bar\gamma+ty)=t^{-3}\rho(y),\qquad
\rho(y):=\frac{1}{\pi}\frac{3y^2-1}{(1+y^2)^3}.
\]
于是
\[
\|r_t\|_{L^p}=t^{1/p-3}\|\rho\|_{L^p}\ \ (p<\infty),\qquad
\|r_t\|_{L^\infty}=t^{-3}\|\rho\|_{L^\infty}.
\]
误差项由同一命题的余项控制给出
\(
\|\varepsilon_t\|_{L^p}=o(t^{1/p-3})
\)，故
\[
\frac{t^{3-1/p}}{\sigma^2}\,\|\tilde h_t-\tilde g_t\|_{L^p}
=t^{3-1/p}\|r_t\|_{L^p}+o(1)\to C_p.
\]
当 \(p=2\) 时，
\(
\int_{\RR}\rho(y)^2\dd y=3/(16\pi)
\)，得到 \eqref{eq:xi-poisson-lp-c2-refined}。
\end{proof}

\begin{theorem}[差分剖面的双零刚性与三阶矩偏心漂移]\label{thm:xi-poisson-difference-two-zero-rigidity-refined}
设 \(\sigma^2=m_2>0\)，并假设
\(
\int_{\RR}\abs{\gamma-\bar\gamma}^4\,\dd\tilde\nu(\gamma)<\infty
\)。
则存在 \(T_0<\infty\)，使得对所有 \(t\ge T_0\)，函数
\(\tilde h_t-\tilde g_t\) 在 \(\RR\) 上恰有两个简单零点
\(
x_-(t)<x_+(t)
\)，并满足
\begin{equation}\label{eq:xi-poisson-difference-two-zero-sign-refined}
\tilde h_t(x)-\tilde g_t(x)<0\ \ (x\in(x_-(t),x_+(t))),\qquad
\tilde h_t(x)-\tilde g_t(x)>0\ \ (x\notin[x_-(t),x_+(t)]).
\end{equation}
若进一步 \(\abs{m_3}<\infty\)，则
\begin{equation}\label{eq:xi-poisson-difference-two-zero-drift-refined}
x_\pm(t)=\bar\gamma\pm\frac{t}{\sqrt3}+\frac{m_3}{3\sigma^2}+o(1)\qquad(t\to\infty).
\end{equation}
\end{theorem}

\begin{proof}
由三阶展开并利用四阶矩控制余项，可得
\[
\tilde h_t=\tilde g_t+\frac{m_2}{2}\tilde g_t''-\frac{m_3}{6}\tilde g_t'''+O(t^{-5})
\]
在 \(x\in\RR\) 上一致成立。令 \(x=\bar\gamma+ty\)，化为
\[
\tilde h_t(\bar\gamma+ty)-\tilde g_t(\bar\gamma+ty)
=\sigma^2t^{-3}f_0(y)+m_3t^{-4}f_1(y)+O(t^{-5}),
\]
其中
\[
f_0(y)=\frac{1}{\pi}\frac{3y^2-1}{(1+y^2)^3},\qquad
f_1(y)=\frac{4}{\pi}\frac{y(y^2-1)}{(1+y^2)^4}.
\]
\(f_0\) 在 \(\pm y_0\)（\(y_0:=1/\sqrt3\)）处有且仅有两个简单零点，且
\(
f_0<0\) 于 \(|y|<y_0\)，\(f_0>0\) 于 \(|y|>y_0\)。
扰动项在 \(t\to\infty\) 下一致趋零，故由简单零点稳定性得到 \eqref{eq:xi-poisson-difference-two-zero-sign-refined}。

令 \(y_\pm(t):=(x_\pm(t)-\bar\gamma)/t\)，设
\(
y_\pm(t)=\pm y_0+c/t+o(t^{-1})
\)。
代入零点方程并线性化：
\[
0=\sigma^2\!\left(f_0'(\pm y_0)\frac{c}{t}+o(t^{-1})\right)
+\frac{m_3}{t}f_1(\pm y_0)+o(t^{-1}).
\]
因此
\[
c=-\frac{m_3f_1(\pm y_0)}{\sigma^2f_0'(\pm y_0)}.
\]
直接计算得 \(-f_1(y_0)/f_0'(y_0)=1/3\)，并由奇偶对称性在 \(-y_0\) 处同值，故 \(c=m_3/(3\sigma^2)\)，即得 \eqref{eq:xi-poisson-difference-two-zero-drift-refined}。
\end{proof}

\begin{proposition}[反向 \texorpdfstring{$D_{\mathrm{KL}}$}{KL} 的 \texorpdfstring{$t^{-6}$}{t^-6} 系数与几何非对称]\label{prop:xi-poisson-reverse-kl-sixth-order-refined}
在定理 \ref{thm:xi-poisson-kl-sixth-order-refined-binomials} 的矩条件下，
\[
D_{\mathrm{KL}}(\tilde g_t\mid \tilde h_t)
:=\int_{\RR}\tilde g_t(x)\log\frac{\tilde g_t(x)}{\tilde h_t(x)}\,\dd x
\]
满足
\begin{equation}\label{eq:xi-poisson-reverse-kl-sixth-order-refined}
D_{\mathrm{KL}}(\tilde g_t\mid \tilde h_t)
=\frac{\sigma^4}{8t^4}
+\frac{1}{t^6}\left(
\frac{\sigma^6}{32}
+\frac{3m_3^2}{32}
-\frac{\sigma^2m_4}{8}
\right)
+o(t^{-6}).
\end{equation}
特别地，
\[
D_{\mathrm{KL}}(\tilde g_t\mid \tilde h_t)-D_{\mathrm{KL}}(\tilde h_t\mid \tilde g_t)
=\frac{\sigma^6}{64t^6}+o(t^{-6}).
\]
\end{proposition}

\begin{proof}
记 \(\delta_t:=\tilde h_t/\tilde g_t-1\)。则
\[
D_{\mathrm{KL}}(\tilde g_t\mid \tilde h_t)
=\int_{\RR}\tilde g_t(x)\,\bigl(-\log(1+\delta_t(x))\bigr)\,\dd x.
\]
由
\(
-\log(1+u)=-u+\frac12u^2-\frac13u^3+O(u^4)
\)，并利用 \(\int\tilde g_t\delta_t=0\)、\(\delta_t=O(t^{-2})\)，有
\[
D_{\mathrm{KL}}(\tilde g_t\mid \tilde h_t)
=\frac12\int_{\RR}G\,\delta_t^2\,\dd y
-\frac13\int_{\RR}G\,\delta_t^3\,\dd y
+o(t^{-6}).
\]
再代入定理 \ref{thm:xi-poisson-kl-sixth-order-refined-binomials} 证明中的 \(\delta_t\) 展开与同一组积分常数，即得
\eqref{eq:xi-poisson-reverse-kl-sixth-order-refined}。
\end{proof}

\begin{corollary}[耗散率 \texorpdfstring{$I_1$}{I1} 的 \texorpdfstring{$t^{-7}$}{t^-7} 精确二项展开]\label{cor:xi-poisson-i1-seventh-order-refined}
定义
\(
I_1(\tilde h_t\mid \tilde g_t):=-\frac{\dd}{\dd t}D_{\mathrm{KL}}(\tilde h_t\mid \tilde g_t)
\)。
在定理 \ref{thm:xi-poisson-kl-sixth-order-refined-binomials} 条件下，
\begin{equation}\label{eq:xi-poisson-i1-seventh-order-refined}
I_1(\tilde h_t\mid \tilde g_t)
=\frac{\sigma^4}{2t^5}
+\frac{1}{t^7}\left(
\frac{3\sigma^6}{32}
+\frac{9m_3^2}{16}
-\frac{3\sigma^2m_4}{4}
\right)
+o(t^{-7}).
\end{equation}
\end{corollary}

\begin{proof}
由 \eqref{eq:xi-poisson-kl-sixth-order-refined-binomials} 逐项求导：
\[
-\frac{\dd}{\dd t}\!\left(\frac{\sigma^4}{8t^4}\right)=\frac{\sigma^4}{2t^5},\qquad
-\frac{\dd}{\dd t}\!\left(\frac{C}{t^6}\right)=\frac{6C}{t^7},
\]
\[
C:=\frac{\sigma^6}{64}+\frac{3m_3^2}{32}-\frac{\sigma^2m_4}{8}.
\]
代入 \(6C\) 并整理即得 \eqref{eq:xi-poisson-i1-seventh-order-refined}。
\end{proof}

\begin{theorem}[Poisson--Cauchy 差分剖面的全形态 \texorpdfstring{$L^1$}{L1} 极限]\label{thm:xi-poisson-cauchy-fullshape-l1-limit}
沿用命题 \ref{con:xi-poisson-coarsegraining-second-order-normal-form} 的设定，
记
\[
\sigma^2:=\int_{\RR}(\gamma-\bar\gamma)^2\,\dd\tilde\nu(\gamma)\in(0,\infty),\qquad
\tilde h_t:=P_t*\tilde\nu,\qquad
\tilde g_t(x):=P_t(x-\bar\gamma),
\]
并定义重标定差分剖面
\[
\Delta_t(y):=\frac{t^3}{\sigma^2}\Bigl(\tilde h_t(\bar\gamma+ty)-\tilde g_t(\bar\gamma+ty)\Bigr),
\qquad y\in\RR.
\]
则当 \(t\to\infty\) 时，
\begin{equation}\label{eq:xi-poisson-cauchy-fullshape-l1-limit}
\Delta_t\longrightarrow R
\quad\text{于 }L^1(\RR,\dd y),
\qquad
R(y):=\frac{1}{\pi}\frac{3y^2-1}{(1+y^2)^3}.
\end{equation}
\end{theorem}

\begin{proof}
由命题 \ref{con:xi-poisson-coarsegraining-second-order-normal-form}，
存在余项 \(\varepsilon_t\) 使
\[
\tilde h_t=\tilde g_t+\sigma^2\tilde r_t+\varepsilon_t,\qquad
\tilde r_t(\bar\gamma+ty)=t^{-3}R(y),\qquad
\|\varepsilon_t\|_{L^1(\RR,\dd x)}=O(t^{-3}).
\]
因此
\[
\Delta_t(y)=R(y)+\frac{t^3}{\sigma^2}\,\varepsilon_t(\bar\gamma+ty).
\]
对误差项作变量代换 \(x=\bar\gamma+ty\)，得到
\[
\|\Delta_t-R\|_{L^1(\RR,\dd y)}
\le
\left\|\frac{t^3}{\sigma^2}\varepsilon_t(\bar\gamma+t\cdot)\right\|_{L^1(\RR,\dd y)}
=\frac{t^2}{\sigma^2}\|\varepsilon_t\|_{L^1(\RR,\dd x)}
=O(t^{-1})\xrightarrow[t\to\infty]{}0.
\]
故 \eqref{eq:xi-poisson-cauchy-fullshape-l1-limit} 成立。证毕。
\end{proof}

\begin{theorem}[Rényi 散度主导系数的普适常数律]\label{thm:xi-poisson-cauchy-renyi-alpha-over-8}
在定理 \ref{thm:xi-poisson-kl-sixth-order-refined-binomials} 的条件下，
对任意固定 \(\alpha\in(0,\infty)\setminus\{1\}\) 定义
\[
D_\alpha(\tilde h_t\mid \tilde g_t)
:=\frac{1}{\alpha-1}\log\!\int_{\RR}\tilde h_t(x)^\alpha\tilde g_t(x)^{1-\alpha}\,\dd x.
\]
则有
\begin{equation}\label{eq:xi-poisson-cauchy-renyi-alpha-over-8}
\lim_{t\to\infty}\frac{t^4}{\sigma^4}\,D_\alpha(\tilde h_t\mid \tilde g_t)=\frac{\alpha}{8}.
\end{equation}
等价地，
\[
D_\alpha(\tilde h_t\mid \tilde g_t)=\frac{\alpha\sigma^4}{8t^4}+o(t^{-4}).
\]
特别地，\(\alpha\downarrow1\) 回收 KL 主导常数 \(\sigma^4/8\)；而 \(\alpha=\tfrac12\) 时，
\[
-\log\!\int_{\RR}\sqrt{\tilde h_t\tilde g_t}\,\dd x
=\frac{1}{2}D_{1/2}(\tilde h_t\mid \tilde g_t)
=\frac{\sigma^4}{32t^4}+o(t^{-4}).
\]
\end{theorem}

\begin{proof}
记
\[
M_\alpha(t):=\int_{\RR}\tilde h_t^\alpha\tilde g_t^{1-\alpha}\,\dd x.
\]
考虑 \(f_\alpha(u):=\frac{u^\alpha-1}{\alpha-1}\)。
则 \(f_\alpha(1)=0\)、\(f_\alpha''(1)=\alpha\)，并且
\[
D_{f_\alpha}(\tilde h_t\mid \tilde g_t)
=\int_{\RR}\tilde g_t(x)\,f_\alpha\!\left(\frac{\tilde h_t(x)}{\tilde g_t(x)}\right)\dd x
=\frac{M_\alpha(t)-1}{\alpha-1}.
\]
由命题 \ref{prop:xi-cauchy-poisson-fdiv-universal-constant} 得
\[
\frac{M_\alpha(t)-1}{\alpha-1}
=\frac{\alpha\sigma^4}{8t^4}+o(t^{-4}),
\]
即
\[
M_\alpha(t)=1+\frac{\alpha(\alpha-1)\sigma^4}{8t^4}+o(t^{-4}).
\]
由于 \(\log(1+u)=u+O(u^2)\) 且 \(u=O(t^{-4})\)，可得
\[
\log M_\alpha(t)
=\frac{\alpha(\alpha-1)\sigma^4}{8t^4}+o(t^{-4}).
\]
两边除以 \(\alpha-1\) 即得 \eqref{eq:xi-poisson-cauchy-renyi-alpha-over-8}。证毕。
\end{proof}

\begin{corollary}[KL 六阶精化的中心矩记号重写]\label{cor:xi-poisson-cauchy-kl-tminus6-mu34-form}
在定理 \ref{thm:xi-poisson-kl-sixth-order-refined-binomials} 的条件下，若进一步满足
\[
\int_{\RR}\abs{\gamma-\bar\gamma}^6\,\dd\tilde\nu(\gamma)<\infty,
\]
并记
\[
\mu_3:=\int_{\RR}(\gamma-\bar\gamma)^3\,\dd\tilde\nu(\gamma),\qquad
\mu_4:=\int_{\RR}(\gamma-\bar\gamma)^4\,\dd\tilde\nu(\gamma),
\]
则
\begin{equation}\label{eq:xi-poisson-cauchy-kl-tminus6-mu34-form}
D_{\mathrm{KL}}(\tilde h_t\mid \tilde g_t)
=\frac{\sigma^4}{8t^4}
+\frac{1}{t^6}\left(
\frac{\sigma^6}{64}
-\frac{\sigma^2\mu_4}{8}
+\frac{3\mu_3^2}{32}
\right)
+o(t^{-6}).
\end{equation}
\end{corollary}

\begin{proof}
该式即 \eqref{eq:xi-poisson-kl-sixth-order-refined-binomials} 在记号
\(
\mu_3=m_3,\ \mu_4=m_4
\)
下的直接改写；六阶矩条件蕴含原定理所需矩条件。证毕。
\end{proof}

\begin{conjecture}[Jensen 缺陷与 Poisson 熵证书的重参数化对齐]\label{conj:xi-jensen-poisson-defect-alignment-refined}
存在仅依赖文内既定坐标门（Cayley/相位门）与缺陷测度构造的显式重参数化
\(
t=t(\rho)\uparrow\infty\ (\rho\uparrow1)
\)，以及由 \(G_t\) 可审计诱导的概率剖面族 \(\{\tilde h_t\}\)，使
\begin{equation}\label{eq:xi-jensen-poisson-defect-alignment-refined}
D_\zeta(\rho)=0\ \ (\forall \rho\in(0,1))
\Longleftrightarrow
\lim_{\rho\uparrow1}t(\rho)^4\,D_{\mathrm{KL}}(\tilde h_{t(\rho)}\mid \tilde g_{t(\rho)})=0.
\end{equation}
若 RH 不成立，则上式极限严格为正，且其首阶量纲由
\eqref{eq:xi-poisson-kl-sixth-order-refined-binomials}
中的 \(\sigma^4/8\) 提供方差指纹。
\end{conjecture}

\endinput
